\chapter{Why Proof Engineering Matters}
\label{cha:appl-mech-proofs}

Formal verification of a program can improve actual and perceived reliability.
It can help the programmer think about the desired and actual behavior of the program,
perhaps finding and fixing bugs in the process~\citep{murraybp}.
It can make explicit which parts of the system are trusted, and further decrease the burden
of trust as more of the system is verified.

One noteworthy program verification success story is the CompCert~\citep{Leroy:POPL06, Leroy2009} verified optimizing C compiler.
Both the back-end and front-end compilation passes
of CompCert have been verified, ensuring the correctness of their composition~\citep{Kaestner2017}.
CompCert has stood up to the trials of human trust: it has been used, for example, to compile code for safety-critical flight control software~\citep{Frana2011}.
It has also stood up to rigorous testing: while the test generation tool Csmith~\citep{Yang2011} found 
79 bugs in GCC and 202 bugs in LLVM, it was unable to find any bugs in the verified parts of CompCert.
%This demonstrated the empirical reliability of CompCert, which in turn further increased confidence in its reliability.

CompCert, however, was not a simple endeavor: The original development comprised of approximately 35,000 lines of Coq code;
functionality accounted for only 13\% of this, while specifications and proofs accounted for the other 87\%.
This is not unusual for large proof developments. The initial correctness proofs for an operating system (OS) kernel in Isabelle/HOL,
for example, consisted of 480,000 lines of specifications and proofs~\citep{Klein2014micro}.

Much like software engineering theories, techniques, and tools help software engineers deal with the challenges
of programming for scale, so proof engineering helps proof engineers.
To highlight the importance of proof engineering, we start by surveying a sample of domains in which 
proof engineering for program verification has been particularly influential (Section~\ref{sec:app-prog}).
We then discuss proof engineering for other domains (Section~\ref{sec:other-domains}),
as the lessons learned from those domains transfer to program verification as well.
We conclude by describing some examples of practical impact from program verification (Section~\ref{sec:empirical}).
%and empirical evaluations (Section~\ref{sec:empirical}).

%\talia{pending better story \& reorg}

%We start by surveying the main areas in which proof mechanization pays
%off, either due to the intricacy of reasoning that becomes
%error-prone, or in correctness-critical software artifacts, such as
%compilers, operational and distributed systems.

\section{Proof Engineering for Program Verification}
\label{sec:app-prog}

We discuss a sample of the domains in which proof engineering has had a large impact:
certified compilers (Section~\ref{sec:verified-compilers}),
low-level systems software (Section~\ref{sec:systems-software}),
concurrent and distributed systems (Section~\ref{sec:distributed-systems}),
and ceritified solvers and checkers (Section~\ref{sec:solvers}).

\subsection{Certified Compilers}
\label{sec:verified-compilers}

Development of certified compilers is a classic application of proof assistants (Section~\ref{sec:early-history}). 
%This is because, among other things,
%compilers are widely used, have broad applications, and often have uncontroversial correctness specifications.
In spite of their long history, however, certified compilers for practical and widely used programming languages have only started to appear in the last
decade. This is mainly due to the sheer size and complexity of the semantics of these languages---challenges that have necessitated 
developments in proof engineering.

The inertia for those developments came in 2006 with the help of CompCert.
This project was a pivotal moment in the history of program verification.
\cite{Leroy:POPL06} received the POPL test of time award in 2016~\citep{popl-time}, with the release noting
its pivotal role:

\begin{quote}
The paper was (and still is) groundbreaking in that it demonstrates the feasibility of using an interactive theorem prover---specifically, Coq---to both program and formally verify a realistic compiler ... [it] made a convincing case that theorem-proving technology is mature enough to be applied to the full functional verification of realistic systems, and in so doing heralded a new age of ``big verification.''
\end{quote} 

Many projects built on CompCert by, for example, using its C semantics, or by targeting C or any of its intermediate representations.
\cite{Appel2011} developed a program logic based on CompCert's C semantics, allowing Coq users to prove properties of deeply embedded C programs in Coq. These properties then hold for machine code generated by CompCert, via CompCert's main correctness result~\citep{Appel:BOOK14}. \cite{Kaestner2017} described many extensions and enhancements to the basic compilation toolchain of CompCert, e.g., translation validation of the process of \emph{linking} machine code to produce object files and executable files. \cite{Cao2018} presented a CompCert-based C program verification environment called VST-Floyd in Coq based on separation logic, simplifying the process of specifying and verifying properties that hold down to machine code.

Certified compilers now span broad applications:
For example, the \textsc{Cogent} language for system programming and verification is accompanied a certifying compiler~\citep{oconnor2016b} 
in Isabelle/HOL that produces a proof that the generated C code is correct. The compiler uses a refinement framework (Section~\ref{sec:refinement}) to
automate relating the \textsc{Cogent} semantics to the generated code~\citep{Rizkallah16}.
The Standard ML language variant CakeML~\citep{Kumar2014} has a verified compiler in HOL4
with a certified machine-code implementation produced by bootstrapping (applying the compiler to itself).
Using the Vellvm~\citep{zhao2012} framework in Coq, proof engineers can reason about transformations on the LLVM intermediate language representation.

Certified compilers have also covered new ground with respect to modularity and compositionality.
The IMM~\citep{Podkopaev2019} memory model modularizes certified compilation from high-level concurrent programming languages to different hardware models.
The Bedrock~\citep{Chlipala2013} intermediate language and verification environment in Coq for low-level programming
contains a notion of macros that can be reasoned about modularly.
The Pilsner~\citep{Neis2015} certified compiler in Coq from a higher-order ML-like language to machine code
supports programs that can, in contrast to CompCert, be compositionally verified.
Compositional CompCert~\citep{Stewart2015} is a variant of CompCert with a correctness theorem that can be applied compositionally. 

Proof engineering for certified compilers has had applications directly to the languages these compilers are verified in.
For example, there are certified compilers both from HOL4~\citep{Myreen2012} and from Isabelle/HOL~\citep{Hupel2018}
to CakeML; it is possible to produce machine code by composing these compilers with the certified CakeML compiler.
Both CertiCoq~\citep{Anand2017} and \OE uf~\citep{Mullen2018} describe certified compilers for Coq's specification
language Gallina. CertiCoq is an ongoing project to build a certified compiler from Gallina to machine code, using a 
hierarchy of custom intermediate languages. \OE uf presents a certified compiler from a subset of Gallina to assembly code.
Both compilers target CompCert intermediate languages: CertiCoq targets Clight, while \OE uf targets Cminor.
Each of these compilers provide an alternative to untrusted program extraction (Section~\ref{sec:trust-programs}).

\subsection{Low-Level Systems Software}
\label{sec:systems-software}

In addition to compilers, low-level software such as operating systems, file systems, and network stacks are building blocks of large software systems. In turn, such software relies on interfaces to hardware, and on hardware behavior. Through considerable effort, proof engineers have specified and verified important pieces of systems software and their hardware bases.

In pioneering work, \cite{Klein2009,Klein2014micro} developed a small general-purpose OS kernel, called seL4, in the C programming language, with correctness proofs in Isabelle/HOL. The proven properties include correctness of interprocess communication, access-control enforcement, and information-flow noninterference. Many verified extensions to seL4 have been proposed since its inception, e.g., to ensure hard deadlines are met for system calls, which is required for applications in real-time systems~\citep{Sewell2017}.

Verification of OS kernels has brought new developments in proof engineering.
For example, as part of the CertiKOS project, \cite{Gu-al:POPL15} presented a framework in Coq for specifying and verifying abstraction layers, which they used to develop several certified OS kernels. In doing so, they introduced the idea of a deep specification (Section~\ref{sec:des-scale}). 
\cite{Gu-al:OSDI16} designed and verified a concurrent kernel for the x86 architecture with fine-grained locking.

An Instruction Set Architecture (ISA) provides an important interface to hardware for, .e.g., compilers.
\cite{Fox2010} developed formal specifications (semantics) of the ISA for ARMv7 in HOL4. \cite{Armstrong2019} later used a domain-specific language to provide ISA specifications for Isabelle/HOL, HOL4, and Coq for the ARMv8, RISC-V, and CHERI-MIPS architectures.
\cite{Morrisett2012} modeled a subset of the x86 ISA in Coq, and used it to build a verified checker for a sandbox policy.

Security policies of systems software are apt targets for verification due to the importance of them being correct.
Along these lines, \cite{Dam2013} proposed a separation kernel (hypervisor) based on the ARMv7 processor architecture and its formalization by Fox and Myreen, and proved, in HOL4, an information flow security property that ensures OS instances can only communicate via explicit channels. \cite{Guanciale2016} proved memory virtualization security in HOL4 for ARMv7.

OS kernel subsystems and formats are additional formalization targets: \cite{Bishop2006,Bishop2018} defined and validated executable specifications in HOL4 for the TCP/IP stack, and \cite{Kell2016} formalized the ELF binary format in HOL4 for executables used in Unix-like operating systems such as Linux.

Verification has also reached file systems, such as FSCQ~\citep{Chen2015}, a file system with guarantees about crash safety that have been verified in Coq,
and DFSCQ~\citep{Chajed2017}, an efficient crash-safe file system with several verified optimizations.
Using the \textsc{Cogent} language and its certifying compiler (Section~\ref{sec:verified-compilers}),
\cite{Amani16} developed a file system called BilbyFS in Isabelle/HOL with executable code in C; they also implemented and verified the legacy Linux file system ext2. \cite{Ridge2015} developed a specification of POSIX file systems in HOL4, which they tested against real-world file system behavior.

\subsection{Concurrent and Distributed Systems}
\label{sec:distributed-systems}

Concurrent and distributed systems can be difficult to develop, understand, and debug. Researchers have developed many different theories and frameworks for reasoning about such systems in proof assistants.
For example, FCSL~\citep{Sergey-al:PLDI15} is a framework in Coq for reasoning about fine-grained concurrency, based on a shallow embedding of concurrent programs.
%that has been used to prove many implementations of classic concurrent data structures correct. % such as lock-based counters and stacks.
The \textsc{Disel}~\citep{Sergey2017} framework for distributed separation logic in Coq builds on this,
using the shallow embedding approach from FCSL. % it has been used to verify an implementation of a distributed two-phase commit algorithm. 
The formalization is extracted to executable OCaml code and run on real hardware.
Two different frameworks~\citep{zeller2014, gomes2017verifying} in Isabelle/HOL exist for verifying two different models
of CRDTs, replicated datatypes
which provide strong eventual consistency guarantees.

Progress in proof engineering for verification of concurrent and distributed systems has had implications for the formalization of practical 
programming languages and protocols. This is illustrated by the \cite{Jung2017} formalization of the imperative and threaded Rust programming language
using the Iris~\citep{Jung2018} framework for concurrent separation logic in Coq,
and by the \cite{Woos2016} implementation and verification of the key correctness property of the Raft consensus algorithm 
using the Verdi~\citep{Wilcox2015} framework for verification of asynchronous message-passing distributed systems.

% TODO Talia: CRDT stuff

%\cite{Pirlea-Sergey:CPP18}

%Verified shims~\citep{Zach}

% FLP https://www.isa-afp.org/entries/FLP.html

% Disk Paxos https://www.isa-afp.org/entries/DiskPaxos.html

%\subsection{Verified Systems and Libraries}
%\label{sec:verif-distr-syst}
%
%\subsubsection{Embedded Systems}
%
%The CompCert verified compiler has found applications in embedded systems~\cite{Kaestner2018}, which includes flight control software~\cite{Frana2011}.
%
%\subsubsection{Security}
%
%The BoringSSL library, used in the widely deployed Google Chrome Web browser, now includes high-performance cryptographic code verified in Coq~\cite{Erbsen2019}.

% possibilities
%\begin{itemize}
%\item Conference management system with verified document confidentiality~\citep{kanav2014conference}
%\item CryptoHOL and its extensions 
%% https://www.isa-afp.org/entries/CryptHOL.html, https://www.isa-afp.org/entries/Game_Based_Crypto.html,
%% https://www.isa-afp.org/entries/Constructive_Cryptography.html
%\item UPF and its extensions
%% https://www.isa-afp.org/entries/UPF.html, https://www.isa-afp.org/entries/UPF_Firewall.html
%\item PikeOS
%% https://www.isa-afp.org/entries/CISC-Kernel.html
%\item For a fun one, hotel key cards
%% https://www.isa-afp.org/entries/HotelKeyCards.html
%\item CertiCrypt
%% http://certicrypt.gforge.inria.fr/
%\end{itemize}

%\subsubsection{Verified Hardware}
%\label{sec:verified-hardware}
%
%\begin{itemize}
%\item \cite{Braibant-Chlipala:CAV13}
%\item SPARCv8 % https://www.isa-afp.org/entries/SPARCv8.html
%\end{itemize}

%\subsubsection{Artificial Intelligence and Machine Learning}
%\label{sec:verif-syst-libr}
%
%\begin{itemize}
%
%\item Characterizing the expressiveness of Deep
%  Learning~\cite{Bentkamp-al:ITP17} % https://www.isa-afp.org/entries/Deep_Learning.html
%
%\end{itemize}

\subsection{Certified Solvers and Checkers}
\label{sec:solvers}

Proof engineers have formalized and proven correct automated solvers for first-order logic and other more restricted logics.
\cite{blanchette2018verified} verified a SAT solver in Isabelle/HOL with conventional features such as clause learning. \cite{schlichtkrull2019verified} verified a purely functional first-order superposition-based solver in Isabelle/HOL and obtained executable code in Standard ML.

Another line of work in Isabelle/HOL formalizes various model checkers, for example
for Linear Temporal Logic~\citep{esparza2013fully} and for timed automata~\citep{wimmer2018verified}.

%IsaFoL https://bitbucket.org/isafol/isafol/wiki/Home

\section{Proof Engineering for Other Domains}
\label{sec:other-domains}

While this survey focuses on proof engineering for program verification,
domains outside of program verification encounter similar proof engineering challenges.
The solutions that these communities develop have implications for proof engineering
for program verification. We discuss these impliciations for two domains:
mathematics (Section~\ref{sec:dsmathematics})
and programming languages metatheory (Section~\ref{sec:dsmetatheory}).

\subsection{Mathematics}
\label{sec:dsmathematics}

% TODO Talia: Look through https://www.isa-afp.org/topics.html and add as relevant throughout mathematics section

Mathematics is a natural application domain for proof assistants.
\textit{Formalized mathematics} is the attempt to formalize mathematical 
theories in part or in whole using proof assistants, so that proofs can be 
mechanically checked. 
Formalized mathematics was one of the first major application domains for proof assistants;
several early ITPs and their predecessors
were designed with this domain in mind (Section~\ref{sec:early-history}).

Early formal developments in mathematics arose in the 
1990s~\citep{zucker1994formalization, van1994checking, bancerek1990fundamental, qed-manifesto-boyer-cade94}.
% something about some ITPs focusing on that
Since then, there have been many notable developments for formalized mathematics,
including the Four Color Theorem~\citep{Gonthier2008}, the Kepler
conjecture~\citep{Hales2011, Hales2017}, the fundamental
theorem of algebra~\citep{Geuvers2000}, G{\"o}del-Rosser incompleteness~\citep{OConnor2005}, and the Jorden curve theorem~\citep{Hales2007}. % TODO and more~\revisit{(cite more)}.
Other interesting mathematical proofs can be found in a comparative overview of different proof assistants for mathematics~\citep{Wiedijk2006}.
The QED manfesto~\citep{qed-manifesto-boyer-cade94} called for a complete database of formalized mathematics.
The UniMath~\citep{UniMath} library is an ongoing attempt to formalize foundations of mathematics~\citep{Voevodsky2015}
in Coq, using homotopy type theory (Section~\ref{sec:equality}).

This section samples tools and design principles for formalized mathematics,
and discusses how they are relevant to program verification.

\paragraph{Design Principles} Many proof developments in mathematics are mature and
involve a large community of contributors. Several of these developments have 
style guides for contributors.
For example, UniMath library has a style guide
that serves to make proofs rigorous, easy to
port to other proof assistants, and less fragile, and to standardize
and improve appearance and readability. Among other things, the style
guide prohibits the addition of new axioms, and encourages the use of
tactics whose semantics are well-defined.  Similarly, the HoTT
library~\citep{Bauer2017} for homotopy type theory contains a style
guide which, among other things, encourages uniform naming principles, outlines methods for
defining equivalences, describes how to use axioms uniformly,
and encourages the use of tactics that have well-defined relationships
with the terms they produce. 

In addition to style guides, proof engineers have developed design principles to handle certain kinds of problems
common in mathematics. \cite{Gonthier2013}, for example, outlines
a number of techniques used in the proof of the Odd Order Theorem.

\cite{Wiedijk2006} compares the styles
of seventeen different proof assistants for mathematics. The book is a collection of
proofs of the irrationality of $\sqrt{2}$ from users of each of the proof assistants,
and a discussion of each of the proof assistants and the proofs in those proof assistants.
This comparison can be useful for understanding the tradeoffs of and design considerations in each of the proof assistants.

\paragraph{Tooling} Formalized mathematics has also seen the development of tooling
to support entire classes of proofs. Notable examples include autarktic
computations for algebraic reasoning~\citep{Barendregt2002}, special support
for equational reasoning in proof checkers~\citep{Barthe1996}, 
decision procedures for fragments of arithemic (Section~\ref{sec:autotactics}),
techniques for reasoning modulo associativity and commutativity~\citep{Braibant2011},
transport methods (Sections~\ref{sec:languagereuse} and~\ref{sec:toolingreuse}),
%the Isabelle rewrite system~\revisit{\citep{nipkow1989term}},
and theory exploration (Section~\ref{sec:autotactics}).

\paragraph{Beyond Mathematics} In formalized
mathematics, user communities of certain frameworks or
libraries adhere to style guides and design principles. These style
guides and design principles have different emphases. 
%For example, the \ssreflect community emphasizes using 
%the small-scale reflection proof methodology (Section~\ref{sec:structuredprooflangs}),
%while both UniMath and the HoTT library emphasize tactics with well-defined semantics
%that are easy to understand.
Proof engineers in communities outside of mathematics may similarly benefit from standardizing
some elements of style depending on the desired outcome.  On a
project-by-project basis, this may make collaboration between proof
engineers easier, and limit the accidental introduction of untrusted
code.

Style guides and design principles also have implications for proof understanding. %In
%general, tactic behavior in languages like Coq is ill-defined, which
%can make proofs difficult for humans to understand. 
Mathematicians are the original proof checkers, so it's perhaps unsurprising that many
mathematics communities emphasize proof understanding by humans. 
Beyond mathematics, human understanding of
proofs communicates information to the reader beyond the theorem
statement itself. Furthermore, just like in software engineering,
effective collaboration between proof engineers hinges on mutual
understanding of the underlying code.

Standardizing style may also make automation easier. For example, by limiting
the set of tactics used within a community, the produced proofs are
more clearly defined, which may make higher-level automation such as
refactoring and repair tools (Section~\ref{sec:evolve}) less challenging.

Proof engineers in communities outside of mathematics may also benefit
from comparitive studies across different proof assistants, similar to
\cite{Wiedijk2006} but for other domains.

Finally, much of the tooling developed for mathematics addresses problems
that occur in proof developments outside of mathematics. 
For example, dealing with equivalences and isomorphisms is a problem that
is not exclusive to the domain of mathematics.
Many of the techniques and tooling from mathematics that solve
this problem may be useful for proof engineers who encounter this 
same problem in other domains.

\subsection{Programming Language Metatheory}
\label{sec:dsmetatheory}

% TODO Talia: Look through https://www.isa-afp.org/topics.html and add as relevant throughout metatheory section

One large domain of focus is \textit{mechanized metatheory}:
proofs about programming languages. The desire to mechanize theory led to the introduction of the 
Edinburgh Logical Framework (LF) by \cite{Harper1987} (and later in more detail, \cite{Harper1993}),
building on ideas from Automath.
LF defined a methodology for encoding and reasoning about a simpler programming language from within the higher-order dependently typed
lambda calculus. Mechanizations of metatheory followed shortly after, both in Nuprl (for example, \cite{howe1988computational})
and in LF (see \cite{Harper2007} for an overview of mechanized metatheory in LF, and \cite{harrison-reflection} for
an early history of mechanized metatheory more broadly).

% TODO Javalight, etc
Since then, the domain has grown to reach practical languages:
The mechanization~ \citep{mechanized-sml} of Standard ML in Twelf formalized
the metatheory of a practical language in its entirety.
WebAssembly has had a formal semantics from the very beginning, which has been mechanized~\citep{watt2018mechanising}
in Isabelle/HOL.
Simplified languages representing the core underlying theories of Scala~\citep{Rompf2016, Amin2017} and of
OCaml~\citep{owens2008} have been verified. 
Results from mechanized metatheory have also influenced verification of real compilers,
like CakeML (Section~\ref{sec:verified-compilers}).

The success of metatheory has brought with it benchmark suites and design principles.
%There have been hundreds of papers on design principles for mechanized metatheory alone.
In addition, mechanized metatheory has influenced new additions to the core languages of ITPs.
This section describes a small sample of those benchmark suites, design principles, and language features,
and discusses how the lessons learned from mechanized metatheory generalize beyond
this domain.

\paragraph{Benchmark Suites} Some of the success of mechanized metatheory is attributable
to benchmark suites that have clearly established the importance of the domain
and set out to define how to measure progress within it. The
\poplmark challenge~\citep{Aydemir2005} has been particularly influential.

The benchmarks in the \poplmark challenge are proofs of properties of
the language System F-Sub~\citep{Cardelli1994}, which has parametric polymorphism and
subtyping. \poplmark highlights specific problems in
proof engineering for metatheory, %: \textit{binding},
%\textit{complex inductions}, \textit{experimentation} with the
%formalization, and \textit{component reuse}. 
and outlines criteria for evaluating the success of technology that
addresses these problems.

%The authors identify
%three criteria for evaluating the success of technology that address
%these problems:

%\begin{enumerate}
%\item It should have \textit{reasonable overhead} compared to proofs by hand.
%\item It should be \textit{transparent} to someone with technical knowledge on formalizing metatheory.
%\item It should have a \textit{reasonable cost of entry} to someone with knowledge in programming languages who is not an expert in using proof assistants.
%\end{enumerate}

15 solutions to the \poplmark challenge remain accessible online~\citep{poplmark-website}. 
Of these solutions, 8 are in Coq, 2 are in Isabelle/HOL,
and the remaining 5 are spread across 5 other ITPs. The solutions cover 8 different ways to represent variable \textit{binders},
one of the problems that \poplmark highlights.

Personal communications~\citep{harpersonal, piercenal2}
suggest that the solution using Twelf~\citep{Pfenning1999} (an LF implementation) was the first solution to solve
all of the difficult parts of the challenge.
The website notes that this solution demonstrated the benefits
of using that framework, including the style of binders it supports, while also sparking an interesting discussion
on different ways of specifying the problem across different frameworks.
Only Arthur Chargu{\'e}raud attempted the same solution in the same proof assistant with different styles of binders;
while his solutions were inconclusive, they inspired later work on making binders easier to represent~\citep{Aydemir2008, Chargueraud2011}.

The \poplmark solutions may be thought of as a springboard for later work. They provided information about what
the state-of-the-art was at the time, which enabled later researchers to measure progress.
Over 300 papers have cited \poplmark since its introduction in 2005.

Still, there is some dissatisfaction with the outcomes of \poplmark.
For example, mst papers that cite \poplmark focus on the difficulty of dealing with
binders, which is just one challenge that proof engineers face when
mechanizing metatheory~\citep{piercenal}. The List-machine Benchmark~\citep{appel2012list},
developed in parallel with \poplmark, deemphasizes
binders and instead emphasizes connections between proofs and real
implementations of compilers. 
\poplmark Reloaded~\citep{abel2017poplmark} emphasizes
logical relations proofs. The ORBI~\citep{felty2018benchmarks, felty2015benchmarks} benchmarks
focus on the tradeoffs of design decisions of different systems for mechanizing metatheory,
rather than on different approaches using a given system.

\paragraph{Design Principles} 
Many papers that cite \poplmark focus on difficulties of dealing with binders. Paper proofs
typically use a \textit{named} representation, where variables are
represented by names. These are easy for humans to reason about, but
make it difficult for tools to reason about alpha-equivalence.
Nominal logics~\citep{Aydemir2006, Urban2008} encode names such that
equivalent terms are alpha-equivalent. These representations have
the benefits of named representations, but without the cost of
difficulty reasoning about alpha-equivalence. While the nominal
approach is common in Isabelle~\citep{Urban2008}, only preliminary
approaches exist in Coq~\citep{Aydemir2006}. %an analogous tool in Coq
%would address many of the difficulties Coq developers currently face
%in mechanizing metatheory. 
Developing such a tool in Coq may be difficult due
to the differences in logics between Coq and Isabelle/HOL, as Nominal Isabelle
makes use~\citep{urban2011} of quotient types~\citep{Homeier05} in HOL. 
% TODO what are the implications of this wrt HOL vs. Coq? Do they imply funext? What about Cyrill's framework for quotients?
% plus, is this possible without quotients?

In the absence of a nominal tool for Coq, proof engineers may
explore tradeoffs between \textit{de Bruijn indexes}
(introduced by \cite{debruijn1972} and explored by \cite{owens2008} and
\cite{schafer2015}, among others) and a \textit{locally nameless}
representation (introduced by \cite{Gordon1993b} and explored by \cite{Leroy2007}, \cite{Aydemir2008},
and \cite{Chargueraud2011}, among others). De Bruijn indexes make reasoning
about alpha-equivalence simple because alpha-equivalence is
definitional equality (Section~\ref{sec:equality}), but they require shifting operations in the code,
which complicates human understanding; 
\cite{Berghofer2007} detail the tradeoffs between de Bruijn indices and names.
Locally nameless attempts to capture the best of both worlds:
it uses names for free variables and de Bruijn indexes
for bound variables. Consequentially, locally nameless requires
reasoning about indexes and shifting operations only for locally closed terms.

LF-based systems like Twelf~\citep{Pfenning1999}, Delphin~\citep{poswolsky2009system}, Beluga~\citep{pientka2008programming} use \textit{higher-order abstract syntax} (HOAS)~\citep{Pfenning1988} to simplify reasoning
about binders. HOAS gives an encoding of binders in the object language (the language that is reasoned about) as binders in the meta-language
(the language of reasoning---in the case of LF, the higher-order dependently typed lambda calculus). % TODO explain meta-language, link to LF somehow
Beluga uses ideas from contextual modal type theory~\citep{nanevski2008contextual} to further simplify reasoning about HOAS encodings.
The Hybrid~\citep{ambler2002} tool makes it possible to use HOAS within Isabelle/HOL;
\cite{capretta2007hybrid} describe a version of Hybrid for Coq, and 
\cite{felty2012hybrid} describe how to use Hybrid to reason about an object language in a manner
similar to Twelf. \cite{felty2010hoas} compare Twelf, Beluga, and Hybrid on case studies
of metatheory using HOAS, and presents
a set of challenge problems which highlights the differences between these systems.
Variants of HOAS such as weak HOAS~\citep{Ciaffaglione2012}
and \textit{parametric higher-order abstract syntax} (PHOAS)~\citep{Chlipala2008}
make HOAS more tractable in general-purpose ITPs like Coq.
%Other approaches include a monadic approach that is largely
%folk knowledge~\revisit{(cite)}, nested abstract syntax~\revisit{(cite)},
%and more~\revisit{(cite)}.

There is a small amount of work on design principles that address the concerns of \poplmark beyond dealing with binders.
For example, Engineering Formal Metatheory~\citep{Aydemir2008} 
identifies specific lemmas that are useful and discusses the organization of theorems, proofs, and automation.
It also introduces \textit{cofinite quantification} of free variables in inductive
relations---defining relations that hold on all but finitely many variables, 
rather than for some fresh variable. This strengthens the premise of the relations,
which in turn strengthens inductive hypotheses for proofs. 

Design principles for mechanized metatheory often go hand-in-hand with high-level frameworks such as 3MT (Section~\ref{sec:metaframe}),
or with domain-specific languages such as Ott and Lem (Section~\ref{sec:specenc}).
Other work in design principles for mechanized metatheory includes an overview
of different ways of formalizing language semantics in an ITP
for the same language~\citep{Bertot2009}, and the use of the coinductive partiality monad~\citep{Capretta2005}
in Agda to define denotational semantics~\citep{Danielsson2012}.
% There is also some work addressing experimentation, which we describe in Section ??. 
% TODO include somewhere: Executable specifications, nondeterminism, mysteries of dropbox, Peter Sewell's work.

% Talia: TODO only if time
%\paragraph{Language Features} The challenges of mechanized metatheory have influenced not just
%design principles and frameworks, but also extensions to the core languages of proof assistants.
%One particularly noteworthy challenge in this regard is that of formalizing type theory
%within type theory~\citep{Dybjer1995, Altenkirch2016}. % TODO add a lot more here
%%Work in this space is a motivating factor for the use of induction-recursion,
%which is supported in both Agda~\revisit{(cite)} and Idris~\revisit{(cite)},  % TODO Talia: Any evidence that they influenced this addition?
%and for induction-induction~\revisit{\citep{forsberg2010inductive}}, % TODO add Jasper's when it comes out, and maybe revisit this citation
%which has some support in Agda~\revisit{(cite)}. % TODO where else?

% TODO John Leo feedback, continued:
% The specific case of Type Theory in Type Theory (where type theory acts as both the proof langauage and object language) has been a focal point of many 
% researchers who use Agda. In fact I believe (I could be wrong) that the addition of induction-recursion and induction-induction to Agda were primarily 
% motivated by the need to strengthen type theory in order to use it to describe itself.
%
%Thorston Altenkirch and his students (Abel, Kaposi, Champman, Kraus, Swierstra, etc) have been particularly acrive in this area.
%There is extensive literature, but a few places to start (and look at their references):
% http://www.cse.chalmers.se/~peterd/papers/inductive.html
% https://akaposi.github.io/#publications
% https://www.ioc.ee/~james/papers/lfmtp08_jmc.pdf
%  https://personal.cis.strath.ac.uk/fredrik.nordvall-forsberg/thesis/thesis.pdf
% Talia: Also http://www.cs.nott.ac.uk/~psztxa/publ/tt-in-tt.pdf


\paragraph{Beyond Metatheory} Few domains have seen as much movement
in the development of design principles for proof engineering as
mechanized metatheory. % Reflecting on the last decade of work in mechanized metatheory can help drive work in other domains.
Opinions on the role that \poplmark played in this are mixed~\citep{piercenal, personappel}. There is little
disagreement that \poplmark was timely: Proof assistants were becoming
more usable, and the ongoing development of CompCert (Section~\ref{sec:verified-compilers})
inspired confidence in their usefulness. At the same time, the properties that researchers wanted to
prove about their languages were becoming larger and more complex. It
was becoming difficult to know that these properties were actually correct,
and to maintain confidence in correctness in the face of changes. 

\poplmark gave a common platform
for experimentation and offered a concrete criteria for success in a
timely domain.  The benchmarks were difficult enough to stress
technology, but simple enough that they were easy to understand and
that experts could prove them in a few weeks.
The work in metatheory that followed \poplmark demonstrated
that general-purpose proof assistants really were usable to prove these properties
that researchers cared about.

Work to this
day continues to use \poplmark as an evaluation metric. 
\cite{Amin2017}, for example, introduce a
proof technique using definitional interpreters that addresses the
open challenge of scaling type soundness proofs to realistic
languages, and evaluate the success of
this technique using the F-Sub language from the \poplmark
benchmarks. This highlights that sometimes, proving
a slightly different property and then showing how that relates to the
original property can be much simpler than proving the original
property directly.

\poplmark suggests that timely benchmark suites are instrumental in
bringing the challenges of design for proof engineering to the
attention of the research community; in doing so, however, they can
narrow the focus to one particularly difficult problem, sometimes to
the exclusion of the bigger picture. Domains outside of metatheory
can take this into consideration.

In addition, the success of experts using LF and Twelf on the \poplmark benchmarks and in mechanizing a practical programming language
suggests that it is worth weighing carefully the tradeoffs of using different ITPs, including ITPs with
special support for a given domain. Along those lines, \cite{Miller2018} argues
that handling of variable bindings should be built into ITPs. 
It is also worth considering the barriers to adoption by non-experts of tools with which experts
have demonstrated success within a domain, and how to overcome those barriers. 
SASyLF~\citep{Aldrich2008}, for example, is one attempt to make LF-based ITPs
more accessible to students.

\section{Practical Impact}
\label{sec:empirical}

Proof engineering has already had a large impact on program verification in many domains,
including those from Section~\ref{sec:app-prog}.
Proof engineers have in recent years verified
operating system~\citep{Klein2009} and web browser~\citep{Jang2012} kernels,
machine learning systems~\citep{DBLP:journals/corr/SelsamLD17},
distributed systems~\citep{Woos2016},
quantum circuits~\citep{rand2017},
constraint solvers~\citep{blanchette2018verified, schlichtkrull2019verified},
compilers~\citep{Leroy:POPL06, Kumar2014},
and file systems~\citep{Chen2015, Amani16}.

So far, proof assistants have had the strongest practical impact in systems software.
The CompCert verified compiler, sold as a commercial product, is finding applications in embedded systems, such as those used in aviation~\citep{CompCert-ERTS-2018}. The BoringSSL library, used in the popular Google Chrome Web browser, recently started to include high-performance cryptographic code in C verified in Coq~\citep{Erbsen2019}. The seL4 verified operating system kernel is used in SCADA systems, and aviation and automotive systems~\citep{Klein2018}.

